%% LaRT AMR Technical Description
%% Author: Kwang-Il Seon
\documentclass[english,a4paper]{kiseon_memo}
\synctex=1
\usepackage{listings}
\usepackage{xcolor}

\lstset{
  basicstyle=\ttfamily\small,
  keywordstyle=\color{blue}\bfseries,
  commentstyle=\color{gray},
  stringstyle=\color{olive},
  frame=single,
  breaklines=true,
  numbers=left,
  numberstyle=\tiny\color{gray},
  language=[95]Fortran
}

\usepackage{babel}
\begin{document}

\title{AMR Extension of LaRT: Octree Structure and Raytrace Algorithm}
\author{Kwang-Il Seon}
\date{2026-04-18}
\maketitle

\tableofcontents
\bigskip

%=============================================================================
\section{Overview}
%=============================================================================

LaRT (Lyman-alpha Radiative Transfer) is a Monte Carlo radiative transfer code
for Lyman-$\alpha$ photon propagation in astrophysical media.
The base code (\texttt{LaRT\_combined}) uses a uniform Cartesian grid.
\texttt{LaRT\_v2.00} extends it with an Adaptive Mesh Refinement (AMR) octree
grid, enabling direct use of outputs from cosmological simulation codes such as
RAMSES.

AMR mode is activated with \texttt{par\%use\_amr\_grid = .true.} in the input
namelist. All photon-physics modules (scattering, raytrace, peel-off, output)
are selected through Fortran procedure pointers set in \texttt{setup.f90};
the outer simulation loop (\texttt{run\_simulation\_mod}) and photon
injection (\texttt{gen\_photon\_car}) are therefore unchanged.

The new source files are:
\begin{center}
\begin{tabular}{ll}
\hline
File & Role \\
\hline
\texttt{octree\_mod.f90}      & Octree data type, tree operations, cell-crossing routines \\
\texttt{grid\_mod\_amr.f90}   & Grid creation, opacity normalisation, frequency grid \\
\texttt{raytrace\_amr.f90}    & Photon propagation through AMR cells \\
\texttt{scattering\_amr.f90}  & Resonance and dust scattering (per-cell parameters) \\
\texttt{peelingoff\_amr.f90}  & Peel-off observer calculation \\
\texttt{read\_ramses\_amr.f90}& Reader for RAMSES binary and generic text AMR formats \\
\hline
\end{tabular}
\end{center}

%=============================================================================
\section{AMR Grid Data Structure}
%=============================================================================

\subsection{Octree Topology}

The AMR grid is represented as a \emph{linear octree}: a flat Fortran array of
cells indexed $1, 2, \ldots, N_\mathrm{cells}$.  Cell~1 is the root, which
covers the entire computational box.  Each internal cell has exactly eight
children obtained by bisecting the parent along all three coordinate axes.

The child octant index is encoded as
\begin{equation}
  i_\mathrm{oct} = 1 + i_x + 2\,i_y + 4\,i_z,
\end{equation}
where $i_x = 1$ if $x \ge c_x$ (the cell centre), otherwise $i_x = 0$, and
similarly for $i_y$ and $i_z$.  Hence octant indices run from 1 to 8.

\subsection{Arrays Stored per Cell}

The Fortran derived type \texttt{amr\_grid\_type} (defined in
\texttt{octree\_mod.f90}) stores the following arrays of size $N_\mathrm{cells}$
for the tree topology:

\begin{center}
\begin{tabular}{lll}
\hline
Array & Dimension & Content \\
\hline
\texttt{parent(:)}      & $N_\mathrm{cells}$        & parent cell index (0 for root) \\
\texttt{children(:,:)}  & $8 \times N_\mathrm{cells}$ & child indices; 0 = leaf \\
\texttt{level(:)}       & $N_\mathrm{cells}$        & refinement level (0 = root) \\
\texttt{cx, cy, cz}     & $N_\mathrm{cells}$        & cell centre coordinates (code units) \\
\texttt{ch(:)}          & $N_\mathrm{cells}$        & cell half-width (code units) \\
\texttt{ileaf(:)}       & $N_\mathrm{cells}$        & $> 0$: leaf index; 0: internal \\
\texttt{icell\_of\_leaf(:)} & $N_\mathrm{leaf}$     & cell index for each leaf \\
\texttt{neighbor(:,:)}  & $6 \times N_\mathrm{cells}$ & face-neighbour table (see Section~\ref{sec:neighbor}) \\
\hline
\end{tabular}
\end{center}

Leaf cells ($\texttt{ileaf}(i) > 0$) carry physical data stored in separate
arrays of size $N_\mathrm{leaf}$:

\begin{center}
\begin{tabular}{lll}
\hline
Array & Symbol & Physical meaning \\
\hline
\texttt{rhokap}  & $\bar\kappa_0$ & HI line-centre opacity per code-unit length \\
\texttt{rhokapD} & $\bar\kappa_d$ & dust opacity per code-unit length \\
\texttt{Dfreq}   & $\Delta\nu_D$ & local Doppler frequency [Hz] \\
\texttt{voigt\_a}& $a$           & Voigt damping parameter $a = A_{21}/(4\pi\Delta\nu_D)$ \\
\texttt{vfx, vfy, vfz} & $\tilde{v}_x,\tilde{v}_y,\tilde{v}_z$ & bulk velocity / local $v_\mathrm{th}$ \\
\hline
\end{tabular}
\end{center}

The spectral output arrays (\texttt{Jout}, \texttt{Jin}, \texttt{Jabs}) have
size $N_\nu$ (the number of frequency bins) and are also members of
\texttt{amr\_grid\_type}.

\subsection{Coordinate Convention}

All cell positions and path lengths are stored in \emph{code units} (kpc, pc, or
whatever unit the input data use).  The conversion factor
\texttt{par\%distance2cm} is used once when computing opacities:
\begin{equation}
  \bar\kappa_0 \;[\mathrm{code\text{-}unit}^{-1}]
    = n_\mathrm{HI}\,[\mathrm{cm}^{-3}]
      \times \frac{\sigma_0}{\Delta\nu_D}
      \times d_\mathrm{cm},
\end{equation}
where $d_\mathrm{cm}$ is the number of centimetres per code unit.  After that,
all quantities are dimensionless with respect to code units, exactly mirroring
the Cartesian convention.

%=============================================================================
\section{Building the Octree}
%=============================================================================

\subsection{Tree Construction: \texttt{amr\_build\_tree}}
\label{sec:build_tree}

\texttt{amr\_build\_tree} takes flat arrays of leaf-cell centre coordinates
(\texttt{xleaf}, \texttt{yleaf}, \texttt{zleaf}), AMR levels
(\texttt{leaf\_level}), and box bounds, and constructs the octree by inserting
each leaf cell top-down from the root.

The algorithm is as follows:
\begin{enumerate}
  \item Initialise the root cell (index 1) with half-width $h = L_\mathrm{box}/2$.
  \item For each leaf $\ell = 1, \ldots, N_\mathrm{leaf}$:
  \begin{enumerate}
    \item Start at the root ($\texttt{icell} = 1$).
    \item For each refinement level from 0 to $\ell_\mathrm{leaf} - 1$:
    \begin{enumerate}
      \item Compute the octant index $i_\mathrm{oct}$ using the leaf centre.
      \item If \texttt{children($i_\mathrm{oct}$, icell)} is zero, create a new
            internal node with half-width $h/2$ and centre offset
            $\pm h/2$ from the parent centre.
      \item Descend: $\texttt{icell} \leftarrow \texttt{children}(i_\mathrm{oct}, \texttt{icell})$.
    \end{enumerate}
    \item Mark the reached cell as leaf $\ell$: $\texttt{ileaf(icell)} = \ell$.
  \end{enumerate}
\end{enumerate}

The array is grown dynamically (doubled) if the pre-allocated size is exceeded.

\subsection{Face-Neighbour Precomputation: \texttt{amr\_build\_neighbors}}
\label{sec:neighbor}

After the tree is built, \texttt{amr\_build\_neighbors} pre-computes a
$6 \times N_\mathrm{cells}$ neighbour table.  Entry
$\texttt{neighbor}(f, i)$ stores the cell index of the neighbour of cell $i$
across face $f$, using the face convention:
\begin{equation}
  f = 1{:}+x,\quad 2{:}-x,\quad 3{:}+y,\quad 4{:}-y,\quad 5{:}+z,\quad 6{:}-z.
\end{equation}
A value of 0 means the face is on the domain boundary.

For each cell $i$, the routine probes the point just outside each face
(offset $h \times 10^{-8}$ from the face) and descends the tree to the cell
at level $\le \ell_i$ that contains that probe point:
\begin{equation}
  \texttt{neighbor}(f, i) = \texttt{amr\_find\_cell\_at\_level}
    \!\left(c_i + h_f\hat{e}_f,\; \ell_i\right),
\end{equation}
where $h_f$ is the signed probe offset along face direction $f$ and
$\hat{e}_f$ is the corresponding unit vector.

The result may be a leaf cell at the same level, a leaf at a coarser level
(if the grid is not refined there), or an internal cell at the same level
(if the neighbour side is \emph{more} refined than cell $i$).  The last case
is handled at runtime by \texttt{amr\_next\_leaf} (Section~\ref{sec:next_leaf}).

This table is built once and then used for O(1) boundary crossings throughout
the simulation.

%=============================================================================
\section{Cell-Crossing Algorithms}
%=============================================================================

\subsection{Finding a Leaf: \texttt{amr\_find\_leaf}}

\texttt{amr\_find\_leaf}$(x, y, z)$ descends the octree from the root,
choosing the correct child octant at each level using the octant formula
(Equation~1), until it reaches a cell with $\texttt{ileaf} > 0$.
The result is the leaf index $\ell \in \{1, \ldots, N_\mathrm{leaf}\}$.
The depth is at most $\ell_\mathrm{max}$, the maximum AMR level.

This function is used only at photon injection (once per photon) and as a
fallback if \texttt{photon\%icell\_amr} is not yet set.

\subsection{Cell-Exit Distance: \texttt{amr\_cell\_exit}}

Given a photon at position $(x, y, z)$ inside cell $i$ with direction
$(k_x, k_y, k_z)$, \texttt{amr\_cell\_exit} computes the distance $t_\mathrm{exit}$
to the nearest cell face and identifies which face $f$ is crossed.

For each axis $\alpha \in \{x, y, z\}$ with cell centre $c_\alpha$ and
half-width $h$, the distances to the $\pm$ faces are:
\begin{equation}
  t_\alpha^+ = \frac{c_\alpha + h - \alpha}{k_\alpha}, \quad
  t_\alpha^- = \frac{c_\alpha - h - \alpha}{k_\alpha},
\end{equation}
with $t = \infty$ (implemented as \texttt{hugest}) when the direction component
is zero or points away from the face.  The exit is the minimum over all six
half-space times:
\begin{equation}
  t_\mathrm{exit} = \min(t_x^+,\, t_x^-,\, t_y^+,\, t_y^-,\, t_z^+,\, t_z^-),
\end{equation}
and $f$ is the index of the minimising element (1–6 in the convention above).

\subsection{Next Leaf After a Face Crossing: \texttt{amr\_next\_leaf}}
\label{sec:next_leaf}

After the photon crosses face $f$ of cell $i$ and the position has been updated
to $(x, y, z)$ at the face, \texttt{amr\_next\_leaf}$(i, f, x, y, z)$ returns
the leaf index of the cell the photon enters.  The algorithm is:

\begin{enumerate}
  \item Look up $\texttt{ineigh} = \texttt{neighbor}(f, i)$.  If zero, the photon
        has left the domain; return 0.
  \item If $\texttt{ileaf(ineigh)} > 0$, the neighbour is already a leaf; return
        $\texttt{ileaf(ineigh)}$.
  \item Otherwise the neighbour is an internal cell at a finer level.  Descend:
        \begin{enumerate}
          \item Compute $i_\mathrm{oct}$ from the actual crossing position
                $(x, y, z)$ and the centre of \texttt{ineigh}.
          \item $\texttt{ineigh} \leftarrow \texttt{children}(i_\mathrm{oct}, \texttt{ineigh})$.
          \item Repeat until $\texttt{ileaf(ineigh)} > 0$.
        \end{enumerate}
  \item Return $\texttt{ileaf(ineigh)}$.
\end{enumerate}

Step~1 is O(1); the descent in step~3 traverses at most
$\Delta\ell = \ell_\mathrm{max} - \ell_i$ extra levels, which is typically
one or two.  In practice this makes cell-boundary crossings effectively O(1)
for any AMR grid.

%=============================================================================
\section{Grid Initialisation}
%=============================================================================

\texttt{grid\_create\_amr} in \texttt{grid\_mod\_amr.f90} performs the
following sequence:

\begin{enumerate}
  \item \textbf{Read leaf-cell data.}  MPI rank~0 reads the data file (RAMSES
        binary or generic text) and broadcasts all arrays to the other ranks.
        Each rank therefore holds a full copy of the tree (appropriate for
        moderate grid sizes).

  \item \textbf{Build octree.}  Call \texttt{amr\_build\_tree} and
        \texttt{amr\_build\_neighbors}.

  \item \textbf{Compute per-cell Doppler and Voigt parameters.}  For leaf $\ell$:
        \begin{equation}
          \Delta\nu_{D,\ell} = \frac{v_{\mathrm{th},1}\sqrt{T_\ell}}{\lambda_0},
          \quad
          a_\ell = \frac{A_{21}}{4\pi\,\Delta\nu_{D,\ell}},
        \end{equation}
        where $v_{\mathrm{th},1} = \sqrt{2k_B/m_H}$ is the thermal velocity per
        $\sqrt{T}$.

  \item \textbf{Compute per-cell opacity and velocity.}
        \begin{align}
          \bar\kappa_{0,\ell} &= n_{\mathrm{HI},\ell}\,\frac{\sigma_0}{\Delta\nu_{D,\ell}}\,d_\mathrm{cm}, \\
          \tilde{v}_{\alpha,\ell} &= \frac{v_{\alpha,\ell}[\mathrm{km\,s}^{-1}]}{v_{\mathrm{th},\ell}[\mathrm{km\,s}^{-1}]},
        \end{align}
        where $\sigma_0$ is the hydrogen Lyman-$\alpha$ cross-section at line
        centre and $d_\mathrm{cm}$ is centimetres per code unit.

  \item \textbf{Normalise opacity to \texttt{taumax}.}  See Section~\ref{sec:taunorm}.

  \item \textbf{Set up frequency grid and output arrays.}
\end{enumerate}

\subsection{Input Data Formats}

Two readers are provided.

\paragraph{RAMSES format (\texttt{amr\_type = 'ramses'}).}
Reads the standard RAMSES AMR and hydro binary output files from
\texttt{output\_XXXXX/}.  Unit scales ($\texttt{unit\_l}$, $\texttt{unit\_d}$,
$\texttt{unit\_t}$) are taken from \texttt{info\_XXXXX.txt}.  The routine
returns leaf positions in cm and $n_H$ in cm$^{-3}$.

\paragraph{Generic text format (\texttt{amr\_type = 'generic'}).}
A plain-text file with a header followed by one line per leaf cell:
\begin{lstlisting}[language={}]
# comment lines start with #
BOXLEN  <box_length_in_code_units>
NLEAF   <number_of_leaf_cells>
# x  y  z  level  nH[cm^-3]  T[K]  vx[km/s]  vy[km/s]  vz[km/s]
<x>  <y>  <z>  <level>  <nH>  <T>  <vx>  <vy>  <vz>
...
\end{lstlisting}
The code unit is set by \texttt{par\%distance\_unit} (e.g., \texttt{'kpc'}).

\subsection{Tau Normalisation by Pole Traversal}
\label{sec:taunorm}

The optical depth normalisation convention in \texttt{LaRT\_v2.00} matches the
Cartesian version: \texttt{taumax} is the line-centre optical depth from the
box centre to the $+z$ boundary along the $z$-axis.

After the raw opacities are set, a ray is traced from the box centre in the
$+z$ direction using the same AMR raytrace kernel:
\begin{equation}
  \tau_\mathrm{pole} = \sum_{\ell}
    \bar\kappa_{0,\ell}\,\phi(0, a_\ell)\,\Delta s_\ell,
\end{equation}
where $\phi(x, a)$ is the Voigt profile at $x=0$ and $\Delta s_\ell$ is the
path-length through leaf $\ell$ along the ray.  The normalisation factor is
\begin{equation}
  f_\mathrm{norm} = \frac{\texttt{taumax}}{\tau_\mathrm{pole}},
\end{equation}
and all opacities are scaled: $\bar\kappa_{0,\ell} \leftarrow f_\mathrm{norm}\,\bar\kappa_{0,\ell}$.
The same factor is applied to dust opacity if present.

This pole-traversal approach is exact for any opacity distribution, in contrast
to the volume-average approximation
$\tau \approx \langle\bar\kappa_0\rangle\,L_\mathrm{box}$
that underestimates $\tau_\mathrm{pole}$ for centrally concentrated clouds such
as a sphere.

%=============================================================================
\section{Photon Propagation}
%=============================================================================

\subsection{Main Raytrace Routine: \texttt{raytrace\_to\_tau\_amr}}
\label{sec:raytrace}

This routine propagates a photon through the AMR grid until it accumulates a
target optical depth $\tau_\mathrm{in}$ (sampled as $-\ln\xi$ for a uniform
random $\xi$) or leaves the box.

The photon state is represented by:
\begin{itemize}
  \item $(x, y, z)$ — current absolute position in code units,
  \item $(k_x, k_y, k_z)$ — unit direction vector,
  \item $x_\nu$ — frequency offset in units of the local Doppler width
        (comoving frame of the current cell),
  \item $\ell$ — current leaf index (\texttt{photon\%icell\_amr}).
\end{itemize}

The loop at each step $n$:
\begin{enumerate}
  \item Convert leaf index to cell index:
        $i = \texttt{icell\_of\_leaf}(\ell)$.
  \item Compute the exit distance and exit face:
        $(t_\mathrm{exit}, f) = \texttt{amr\_cell\_exit}(i, x, y, z, k_x, k_y, k_z)$.
  \item Evaluate the opacity at the current frequency:
        \begin{equation}
          \kappa = \bar\kappa_{0,\ell}\,\phi(x_\nu, a_\ell)
                  + \bar\kappa_{d,\ell}
                  \quad (\text{with dust, if present}),
        \end{equation}
        where $\phi(x, a)$ is the Voigt function.
  \item Update accumulated optical depth:
        $\tau \leftarrow \tau + \kappa\,t_\mathrm{exit}$.
  \item \textbf{If $\tau \ge \tau_\mathrm{in}$}: backtrack to the exact
        interaction point (Section~\ref{sec:overshoot}) and exit the loop.
  \item Otherwise: advance position to the face,
        \begin{equation}
          (x, y, z) \leftarrow (x + t_\mathrm{exit}\,k_x,\;
                                 y + t_\mathrm{exit}\,k_y,\;
                                 z + t_\mathrm{exit}\,k_z).
        \end{equation}
  \item Look up the next leaf:
        $\ell' = \texttt{amr\_next\_leaf}(i, f, x, y, z)$.
  \item If $\ell' = 0$, the photon has escaped; set
        \texttt{photon\%inside = .false.} and exit.
  \item Apply the frequency shift (Section~\ref{sec:freqshift}):
        $x_\nu \leftarrow x_\nu'$.
  \item Set $\ell \leftarrow \ell'$ and continue.
\end{enumerate}

\subsection{Frequency Shift at Cell Boundaries}
\label{sec:freqshift}

Each cell has its own Doppler frequency $\Delta\nu_{D,\ell}$ (because cells
can have different temperatures) and its own bulk velocity component along the
photon direction,
\begin{equation}
  u_\ell = \tilde{v}_{x,\ell}\,k_x
          + \tilde{v}_{y,\ell}\,k_y
          + \tilde{v}_{z,\ell}\,k_z
  \quad [\text{dimensionless, in units of local } v_\mathrm{th}].
\end{equation}

When crossing from cell $\ell$ (old) to cell $\ell'$ (new), the photon
frequency must be transformed to the comoving frame of the new cell.  The
physical frequency is invariant, so:
\begin{equation}
  \nu_\mathrm{phys} = \nu_0\left(1 + \frac{(x_\nu^{\rm old} + u_\ell)\,\Delta\nu_{D,\ell}}
                                          {\nu_0}\right),
\end{equation}
and the new dimensionless frequency is:
\begin{equation}
  \boxed{
  x_\nu^{\rm new} = (x_\nu^{\rm old} + u_\ell)
                    \,\frac{\Delta\nu_{D,\ell}}{\Delta\nu_{D,\ell'}}
                    - u_{\ell'}.
  }
\end{equation}
This formula transforms between the comoving frames of two cells with
different temperatures and bulk velocities in a single step.

\subsection{Overshoot Correction}
\label{sec:overshoot}

When $\tau > \tau_\mathrm{in}$ after stepping across a cell, the photon
actually stopped inside the cell.  The exact interaction distance is found
by backtracking:
\begin{equation}
  \Delta s_\mathrm{back} = \frac{\tau - \tau_\mathrm{in}}{\kappa}.
\end{equation}
The total displacement from the entry point is then
$d = d + t_\mathrm{exit} - \Delta s_\mathrm{back}$, and the final position is
\begin{equation}
  (x, y, z) \leftarrow (\texttt{photon\%}x + d\,k_x,\;
                         \texttt{photon\%}y + d\,k_y,\;
                         \texttt{photon\%}z + d\,k_z).
\end{equation}
Note that the position update is accumulated as a single floating-point
addition from the \emph{original} entry point (not from successive face crossings),
which preserves numerical precision.

\subsection{Photon Escape and Spectrum Accumulation}

When the photon leaves the box ($\ell' = 0$), the frequency is converted to
the lab frame:
\begin{equation}
  x_\nu^\mathrm{lab} = x_\nu + u_\ell,
\end{equation}
then scaled to the reference Doppler frequency:
\begin{equation}
  x_\nu^\mathrm{ref} = x_\nu^\mathrm{lab}
                       \,\frac{\Delta\nu_{D,\ell}}{\Delta\nu_{D,\mathrm{ref}}}.
\end{equation}
The frequency bin index is
\begin{equation}
  i_\nu = \left\lfloor
    \frac{x_\nu^\mathrm{ref} - x_\mathrm{min}}{\delta x}
  \right\rfloor + 1,
\end{equation}
and the photon weight is accumulated into $\texttt{amr\_grid\%Jout}(i_\nu)$.

\subsection{Other Raytrace Routines}

Six additional routines in \texttt{raytrace\_amr\_mod} share the same AMR
traversal kernel but accumulate different quantities:

\begin{center}
\begin{tabular}{ll}
\hline
Routine & Purpose \\
\hline
\texttt{raytrace\_to\_edge\_amr}            & Total $\tau$ to box edge (photon unchanged) \\
\texttt{raytrace\_to\_edge\_tau\_gas\_amr}  & Gas-only $\tau$ to box edge \\
\texttt{raytrace\_to\_edge\_column\_amr}    & $N_\mathrm{HI}$ and dust $\tau$ to box edge \\
\texttt{raytrace\_to\_dist\_amr}            & Total $\tau$ to fixed distance $s$ \\
\texttt{raytrace\_to\_dist\_tau\_gas\_amr}  & Gas-only $\tau$ to fixed distance $s$ \\
\texttt{raytrace\_to\_dist\_column\_amr}    & $N_\mathrm{HI}$ and dust $\tau$ to fixed distance $s$ \\
\hline
\end{tabular}
\end{center}

All six routines propagate a copy of the photon state (photon is not modified)
and apply the same frequency-shift formula at each cell crossing so that the
Voigt opacity is evaluated at the correct comoving frequency in every cell.

%=============================================================================
\section{Scattering}
%=============================================================================

The scattering routines in \texttt{scattering\_amr.f90} mirror the Cartesian
versions but read physical parameters from \texttt{amr\_grid} arrays indexed
by the leaf index \texttt{photon\%icell\_amr}.

\subsection{Resonance Scattering}

\texttt{scatter\_resonance\_nostokes\_amr} implements the resonance scattering
algorithm (no Stokes polarisation).  The scattering is handled by the procedure
pointer \texttt{do\_resonance}, which is set at initialisation to one of
\texttt{do\_resonance1\_amr}, \texttt{do\_resonance2\_amr}, etc., depending on
the line type.

\begin{enumerate}
  \item \textbf{Sample the scattering atom velocity.}
        The thermal velocity component of the atom along the photon direction
        (in the cell comoving frame) is sampled from the redistribution
        function via \texttt{rand\_resonance\_vz}:
        \begin{equation}
          u_z = x_\nu + u_\ell.
        \end{equation}
        The two perpendicular thermal velocity components $(u_x, u_y)$ are
        drawn independently from a Gaussian with unit variance.

  \item \textbf{Sample the scattering angle.}
        The polar scattering angle $\theta$ (between the incoming and outgoing
        photon directions) is drawn from the phase function
        \begin{equation}
          P(\cos\theta) = \frac{3}{8}E_1\cos^2\theta + \frac{4 - E_1}{8},
        \end{equation}
        via the function \texttt{rand\_resonance}$(E_1)$.  The parameter $E_1$
        is a property of the atomic transition:
        \begin{itemize}
          \item $E_1 = 1$: pure Rayleigh scattering ($J=0 \to J=1 \to J=0$, e.g.\ He~\textsc{i}).
          \item $0 < E_1 < 1$: mixed Rayleigh + isotropic
                (e.g.\ Ly$\alpha$: $E_1 \approx 0$ for the $J=1/2$ branch,
                $E_1 = 1$ for the $J=3/2$ branch, giving an effective
                $E_1$ that depends on the occupation of upper sub-levels).
          \item $E_1 = 0$: isotropic scattering.
          \item $-1/2 \le E_1 < 0$: oblate phase function (forward and backward
                directions suppressed).
        \end{itemize}
        The azimuthal angle $\phi$ is drawn uniformly from $[0, 2\pi)$.

  \item \textbf{Compute the new photon frequency.}
        In the atom frame the frequency after scattering is
        \begin{equation}
          x'_\nu = x_\nu^\mathrm{atom}
                  + u_z \cos\theta
                  + (u_x\cos\phi + u_y\sin\phi)\sin\theta,
        \end{equation}
        where $x_\nu^\mathrm{atom} = x_\nu - u_z$ is the frequency in the
        atom rest frame.  An optional recoil correction
        $\Delta x = -g_\mathrm{recoil}(1-\cos\theta)$ is applied if
        \texttt{par\%recoil = .true.}

  \item \textbf{Apply core-skip acceleration} if $|x_\nu| < x_\mathrm{crit}$.
\end{enumerate}

The core-skip threshold $x_\mathrm{crit}$ is computed in
\texttt{grid\_create\_amr} from
$a\tau_0^\mathrm{cell}$, the product of the mean Voigt damping parameter and
the mean optical depth per AMR leaf cell.

\subsection{Dust Scattering}

\texttt{scatter\_dust\_nostokes\_amr} uses the standard
Henyey--Greenstein phase function.  When dust is present, the probability of
the interaction being a dust event is
\begin{equation}
  p_d = \frac{\bar\kappa_{d,\ell}}
             {\bar\kappa_{0,\ell}\,\phi(x_\nu, a_\ell) + \bar\kappa_{d,\ell}},
\end{equation}
sampled at the scattering location (cell $\ell$).

\subsection{Multi-Level Resonance Transitions}

For complex atoms (HeI, FeII) represented by multiple transitions, the AMR
scattering module provides \texttt{do\_resonance4\_amr},
\texttt{do\_resonance5\_amr}, and \texttt{do\_resonance6\_amr}, which follow
the same logic as their Cartesian counterparts but use per-cell Doppler
frequencies and Voigt parameters.

%=============================================================================
\section{Output Normalisation}
%=============================================================================

After all photons are traced, the MPI reduction collects all per-frequency
contributions via \texttt{MPI\_ALLREDUCE} into \texttt{amr\_grid\%Jout}
(\texttt{output\_reduce\_amr} in \texttt{output\_sum\_rect.f90}).

The normalised specific intensity (per unit frequency per steradian) is:
\begin{equation}
  J_\mathrm{out}(x_\nu) = \frac{\texttt{Jout}(i_\nu)}
    {N_\mathrm{ph}\;\delta x\;2\pi\;A},
\end{equation}
where $N_\mathrm{ph}$ is the number of photons, $\delta x$ the frequency bin
width, and
\begin{equation}
  A = 4\pi\,d_\mathrm{cm}^2
\end{equation}
is the surface area of a unit sphere in cm$^2$ (the factor $d_\mathrm{cm}$
is \texttt{par\%distance2cm}).

For dimensionless runs (\texttt{distance2cm = 1}, no \texttt{distance\_unit}),
this reduces to $A = 4\pi$, making $J_\mathrm{out}$ independent of the
arbitrary box size.  This matches the Cartesian convention
($A = 4\pi\,r_\mathrm{max}^2\,d_\mathrm{cm}^2$) when
$r_\mathrm{max} = L_\mathrm{box}/2$ and the same \texttt{distance2cm} is used.

%=============================================================================
\section{Core-Skip Acceleration}
%=============================================================================

Lyman-$\alpha$ photons scatter extremely frequently in the line core
($|x_\nu| \ll 1$), where the optical depth per cell is large.  The
core-skip algorithm (Laursen et al. 2009) exploits the fact that a photon in
the core will scatter many times within a single cell without travelling far.
When $|x_\nu| < x_\mathrm{crit}$, the frequency is redrawn from the
wing distribution directly, skipping the core-scattering steps.

In AMR mode, $x_\mathrm{crit}$ is computed from the mean opacity product:
\begin{equation}
  (a\tau_0)_\mathrm{cell}
    = \bar{a}\;\bar\tau_0\;/\;N_\mathrm{cell,axis},
\end{equation}
where $N_\mathrm{cell,axis} = L_\mathrm{box} / (2\,h_\mathrm{root})$ is the
number of root cells along one axis.  The threshold is then
\begin{equation}
  x_\mathrm{crit} =
  \begin{cases}
    0.02\,\exp\!\left[\xi\,(\ln(a\tau_0)_\mathrm{cell})^\chi\right] &
      (a\tau_0)_\mathrm{cell} > 1, \\
    0 & \text{otherwise},
  \end{cases}
\end{equation}
with $(\xi, \chi) = (0.6, 1.2)$ for $(a\tau_0)_\mathrm{cell} \le 60$ and
$(1.4, 0.6)$ otherwise.

%=============================================================================
\section{Validation}
%=============================================================================

The AMR implementation was validated against equivalent Cartesian runs
for a uniform-density sphere of HI gas with a central point source
($T = 10^4\,\mathrm{K}$, no dust).

The AMR sphere is described by the generic text format with the same set of
leaf cells at the same AMR level, providing effectively the same resolution
as the $101^3$ Cartesian grid.  The Cartesian comparison uses
$r_\mathrm{max} = L_\mathrm{box}/2$ with no \texttt{distance\_unit}.

\begin{center}
\begin{tabular}{ccccc}
\hline
$\tau_0$ & $J_\mathrm{out,max}^\mathrm{AMR}$
         & $J_\mathrm{out,max}^\mathrm{CAR}$
         & Ratio (AMR/CAR)
         & Peak velocity \\
\hline
$10^2$   & $6.630 \times 10^{-3}$ & $6.836 \times 10^{-3}$ & 0.970 & $\pm 26.75\,\mathrm{km\,s}^{-1}$ \\
$10^4$   & $7.144 \times 10^{-3}$ & $7.402 \times 10^{-3}$ & 0.965 & $\pm 38.21\,\mathrm{km\,s}^{-1}$ \\
\hline
\end{tabular}
\end{center}

Peak positions, column densities $N(\mathrm{HI})_\mathrm{pole}$, and
$\tau_\mathrm{pole}$ agree exactly between AMR and Cartesian runs.  The
$\sim 3$--$5\%$ amplitude difference is entirely attributable to the
octree voxelisation of the sphere boundary: leaf cells that straddle the
sphere edge carry only the average density, softening the sharp boundary.
No normalisation error is present.

%=============================================================================
\section{Summary of Key Design Choices}
%=============================================================================

\begin{enumerate}
  \item \textbf{Flat linear octree.}  All cells (internal and leaf) are stored
        in a single 1-indexed Fortran array.  Parent--child relationships are
        maintained through integer index arrays.  This avoids pointer
        indirection and is cache-friendly.

  \item \textbf{Precomputed face-neighbour table.}  Built once before the
        simulation loop.  Boundary crossings cost O(1) (plus O($\Delta\ell$)
        descent into finer cells if the neighbour is more refined).

  \item \textbf{Per-cell Voigt parameters.}  Each leaf stores its own
        $\Delta\nu_D$ and $a$, supporting arbitrary temperature distributions
        without any approximation.

  \item \textbf{Frequency shift at every cell crossing.}  The comoving-frame
        frequency is updated at each boundary, so the Voigt opacity is always
        evaluated in the correct local frame, including bulk velocity effects.

  \item \textbf{Single floating-point accumulation of displacement.}  The
        photon position is updated as $\mathbf{r}_0 + d\,\hat{k}$ rather than
        as a sequence of increments, preserving numerical precision across
        many cell crossings.

  \item \textbf{Procedure pointers for AMR/Cartesian switching.}
        The simulation loop and photon injection are unchanged.  AMR physics
        is activated solely by reassigning Fortran procedure pointers in
        \texttt{setup\_v2c.f90}.

  \item \textbf{Pole-traversal tau normalisation.}  The AMR raytrace kernel
        itself is used to compute $\tau_\mathrm{pole}$, giving an exact
        normalisation that is consistent with the Cartesian convention and
        avoids volume-average errors.
\end{enumerate}

\end{document}
